%======================================================================
\section{Proofs}
\label{app:proofs}
%======================================================================

Throughout, $\Lat$ is a periodic lattice of sites, $B_R(j)=\{k\in\Lat:
\dist(k,j)\le R\}$ is the ball of radius $R$ sites around $j$, and for a
lattice field $a$ we write $a|_S$ for its restriction to a set of sites
$S$. All maps act on a batch of fields, but the batch index plays no role
and is suppressed.

\para{Locality.}
A map $G$ between lattice fields has \emph{reach} $R$ if, for every site
$j$, $[G(a)]_j$ depends on $a$ only through $a|_{B_R(j)}$; that is,
$a|_{B_R(j)}=b|_{B_R(j)}$ implies $[G(a)]_j=[G(b)]_j$. A per-site map
(a $1\times1$ convolution, a pointwise nonlinearity, layer normalization
over channels, a per-site MLP) has reach $0$. A circular convolution of
radius $r$ has reach $r$. Reaches add under composition: if $G_1$ has
reach $R_1$ and $G_2$ has reach $R_2$, then $[G_2(G_1(a))]_j$ depends on
$G_1(a)|_{B_{R_2}(j)}$, each entry of which depends on $a$ within $R_1$
of its own site, so $G_2\circ G_1$ has reach at most $R_1+R_2$. A residual
map $a\mapsto a+G(a)$ has the reach of $G$. Hence one block of
Eq.~\eqref{eq:mixer_block} has reach $w$, and the propagator $\Fp$ of
Eq.~\eqref{eq:propagator}, with pointwise embedding and projection around
$n_\ell$ blocks, has reach $\Rprop = w\,n_\ell$. The same count gives
$\Renc$ and $\Rdec$ for the encoder and decoder, measured in sites of
the patch lattice.

%----------------------------------------------------------------------
\subsection{Proof of Proposition~\ref{prop:equivariance}}
%----------------------------------------------------------------------

Each map is a composition of the following operations, so it suffices to
check that each one intertwines the shifts.

\emph{Patch embedding.} With a kernel $K$ supported on one patch and
stride $p$, the embedding is $(\mathcal{P}u)_j =
\sum_{\bm{\xi}\in\text{patch}}K(\bm{\xi})\,u(\bm{x}_j+\bm{\xi})$, where
$\bm{x}_j$ is the corner of patch $j$. Then
\[
\begin{aligned}
  (\mathcal{P}\,\Tr_{s}u)_j
  &= \sum_{\bm{\xi}}K(\bm{\xi})\,u(\bm{x}_j+\bm{\xi}-s h)\\
  &= \sum_{\bm{\xi}}K(\bm{\xi})\,u(\bm{x}_{j-s}+\bm{\xi})\\
  &= (\mathcal{P}u)_{j-s}
  = (\Shift_s\mathcal{P}u)_j,
\end{aligned}
\]
using $\bm{x}_j - sh = \bm{x}_{j-s}$, which holds because the shift is by
a whole number of patches. The gauge anchor of Eq.~\eqref{eq:anchor} is
the special case $K\equiv1/|\text{patch}|$.

\emph{Circular convolution on the lattice.} $(\kappa * a)_j =
\sum_{|\delta|\le r}\kappa_\delta\,a_{j+\delta}$ with indices taken
modulo the lattice size. Then $(\kappa*\Shift_s a)_j = \sum_\delta
\kappa_\delta a_{j+\delta-s} = (\kappa*a)_{j-s} = (\Shift_s(\kappa*a))_j$.

\emph{Per-site operations} ($1\times1$ convolutions, GELU, layer
normalization over channels, per-site MLPs, the residual sum) act
identically at every site and therefore commute with any permutation of
sites, in particular with $\Shift_s$.

\emph{Transposed patch convolution.} The decoder output is
$\Dec(a)(\bm{x}) = \sum_{j}K'(\bm{x}-\bm{x}_j)\,a_j$ with $K'$ supported
on one patch. Then
\[
\begin{aligned}
  \Dec(\Shift_s a)(\bm{x}) &= \sum_j K'(\bm{x}-\bm{x}_j)\,a_{j-s}\\
  &= \sum_k K'(\bm{x}-sh-\bm{x}_k)\,a_k\\
  &= \Dec(a)(\bm{x}-sh) = (\Tr_{s}\Dec(a))(\bm{x}).
\end{aligned}
\]

Composing these identities gives Eq.~\eqref{eq:equivariance}. For the
second claim, every kernel and weight matrix above has a shape fixed by
$(p, c, H, r_e, n_e, r_d, n_d, d_m, w, n_\ell)$ and the MLP ratio; none
depends on $P$. The same $\theta$ therefore defines the maps for every
$P$; the proviso in the statement, at least $2r+1$ sites along every
axis on which a kernel of radius $r$ acts, ensures that each circular
convolution reads $2r+1$ distinct sites. \qed

\begin{remark}
The proviso is not cosmetic. When a lattice axis has fewer than $2w+1$
sites, a circular kernel of radius $w$ overlaps itself, so a site
interacts with copies of itself and its neighbours more than once. The
1D model transferred to $L=25$ has $4$ sites against a propagator kernel
of $7$; it is the one size at which the model does not reproduce the
truth, and its behaviour there depends on the initial condition
(Sec.~\ref{sec:results1d}).
\end{remark}

%----------------------------------------------------------------------
\subsection{Proof of Proposition~\ref{prop:halo}}
%----------------------------------------------------------------------

\emph{Setting.} Let $\Lat$ be any periodic lattice and $C\subset\Lat$ a
crop: a set of contiguous sites along the cropped axes, with the outer
$\Renc$ sites already removed, so that the encoding of the crop agrees
with the restriction of the encoding of the full field on $C$ (this is
the locality statement above applied to $\E$). Let $(z^m)_{m\ge0}$ be the
true latent trajectory on $\Lat$. Define the free interior
\[
  I = \{\, j\in C : \dist(j,k) > \eta \ \text{for every } k\in\Lat\setminus C \,\}
\]
and the halo $H=C\setminus I$. Let $\Fp_C$ denote the propagator applied
to a field on $C$ with \emph{any} fixed rule for the values it reads
outside $C$ (zeros, circular wrap of the crop, or anything else), and
$\Fp_\Lat$ the propagator on $\Lat$.

The forced-halo rollout is
\[
  \hat z^0 = z^0|_C, \quad
  \hat z^{m+1}|_I = [\Fp_C(\hat z^m)]_I, \quad
  \hat z^{m+1}|_H = z^{m+1}|_H .
\]
The reference sequence is the propagator on $\Lat$ with every site
outside $I$ held at its true value:
\[
\begin{gathered}
  \tilde z^0 = z^0, \qquad
  \tilde z^{m+1}|_I = [\Fp_\Lat(\tilde z^m)]_I,\\
  \tilde z^{m+1}|_{\Lat\setminus I} = z^{m+1}|_{\Lat\setminus I}.
\end{gathered}
\]

\emph{Claim.} If $\eta\ge\Rprop$, then $\hat z^m|_I = \tilde z^m|_I$ for
every $m\ge0$.

\emph{Proof.} For $j\in I$ every $k$ with $\dist(k,j)\le\Rprop\le\eta$
lies in $C$, by the definition of $I$. Hence $B_{\Rprop}(j)\subset C$,
the rule $\Fp_C$ uses for sites outside $C$ is never read when computing
$[\Fp_C(a)]_j$, and by locality
\begin{equation}
  a|_C = b|_C \ \Longrightarrow\ [\Fp_C(a)]_j = [\Fp_\Lat(b)]_j
  \qquad (j\in I),
  \label{eq:halo_locality}
\end{equation}
where $b$ is any field on $\Lat$ extending $a$. We show by induction that
$\hat z^m = \tilde z^m|_C$. At $m=0$ both equal $z^0|_C$. If
$\hat z^m=\tilde z^m|_C$, then Eq.~\eqref{eq:halo_locality} with
$a=\hat z^m$, $b=\tilde z^m$ gives $\hat z^{m+1}|_I=\tilde z^{m+1}|_I$,
and on $H\subset\Lat\setminus I$ both sequences equal $z^{m+1}$. Hence
$\hat z^{m+1}=\tilde z^{m+1}|_C$. \qed

\emph{Failure for $\eta<\Rprop$.} Then some $j\in I$ has a site
$k\notin C$ with $\dist(k,j)\le\Rprop$. Unless the weights happen to make
$[\Fp(a)]_j$ independent of $a_k$ (a non-generic case), changing the
true value $z^0_k$ changes $\tilde z^1_j$ but not $\hat z^1_j$, which
never reads $z^0_k$. So the two sequences differ for some truth. The
numerical test reported in Sec.~\ref{sec:crop_training} realizes this
with $\eta=\Rprop-1$.

\begin{remark}
The reference sequence $\tilde z$ is a limited-area model: the
propagator runs freely on $I$ and receives exact lateral boundary data.
On the periodic deployment domain there is no forcing, and sites in $I$
within $\Rprop$ of $H$ see predicted rather than true neighbours after
the first step. The training objective is therefore exact for the
limited-area problem and approximate, in a strip of width $\Rprop$, for
the deployed one.
\end{remark}

%----------------------------------------------------------------------
\subsection{Proof of Lemma~\ref{lem:mse}}
%----------------------------------------------------------------------

Write $m(\tilde z^0)=\Exp[z^k\mid\tilde z^0]$. For any measurable $\Phi$
with $\Exp\|\Phi(\tilde z^0)\|^2<\infty$,
\[
  \Exp\|\Phi(\tilde z^0)-z^k\|^2
  = \Exp\|\Phi(\tilde z^0)-m(\tilde z^0)\|^2
  + \Exp\|m(\tilde z^0)-z^k\|^2,
\]
because the cross term
$2\,\Exp\big[(\Phi-m)^{\top}\Exp[m-z^k\mid\tilde z^0]\big]$ vanishes. The
first term is non-negative and zero for $\Phi=m$, so $\Phi^\star=m$. The
law of total covariance,
$\Cov(z^k)=\Cov\big(\Exp[z^k\mid\tilde z^0]\big)+\Exp\,\Cov(z^k\mid\tilde z^0)$,
gives Eq.~\eqref{eq:total_variance} after taking the trace. \qed

\begin{remark}
If the system is mixing, $\Exp[z^k\mid\tilde z^0]$ tends to the
climatological mean $\Exp[z]$ as $k\to\infty$, so the squared-error-optimal
$k$-step predictor tends to a constant and the first term of
Eq.~\eqref{eq:total_variance} tends to zero. An autonomous propagator
cannot be the conditional mean at every lead time simultaneously, since
its $k$-step map is the $k$-fold iterate of one map. The lemma therefore
describes the direction in which the rollout loss pulls the propagator
(toward iterates with less spread than the data), not the map it
converges to.
\end{remark}

%----------------------------------------------------------------------
\subsection{Proof of Proposition~\ref{prop:blockdiag}}
%----------------------------------------------------------------------

With $\Rdec=0$ the decoder has no lattice convolutions: a $1\times1$
convolution and a pointwise nonlinearity act on each site, and the
transposed convolution with kernel and stride equal to the patch writes
site $j$'s output only onto patch $j$. Hence $\Dec(z)$ restricted to patch
$j$ is $g(z_j)$ for one fixed function $g:\R^c\to\R^{p^{\dim}}$. Let
observation $i$ read grid point $\bm{x}_i$ in patch $s(i)$. Then
$h_i(z)=g(z_{s(i)})(\bm{x}_i)$ does not depend on $z_j$ for $j\ne s(i)$,
which is the block-diagonal structure.

\emph{(a)} Let $v^{(k)}=\sum_{j\in\Lat}e_{j,k}$. Then
\[
  (Jv^{(k)})_i = \sum_{j\in\Lat}\sum_{l=1}^{c} J_{i,(j,l)}\,v^{(k)}_{(j,l)}
  = \sum_{j\in\Lat} J_{i,(j,k)} = J_{i,(s(i),k)},
\]
since all terms with $j\ne s(i)$ vanish. So the $c$ products
$Jv^{(1)},\dots,Jv^{(c)}$ contain every possibly non-zero entry of $J$,
and every other entry is zero. Each product is one forward-mode pass
through $h$, whatever $P$ and $m$ are.

\emph{(b)} For site $s$, let $O_s$ be the observations with non-zero
taper weight and $U_s=s(O_s)$ the sites that own them. The
perturbed-observation analysis of site $s$ uses the rows $O_s$ of three
quantities: the innovations, computed from the nonlinear $h$ evaluated
once for the whole ensemble; the cross-covariance
$\Cov(z_s, z)\,J_{O_s,:}^{\top}$; and the observation-space covariance
$J_{O_s,:}\,B\,J_{O_s,:}^{\top}$. By the block structure, the rows
$J_{O_s,:}$ vanish outside the columns of $U_s$, so
\[
\begin{aligned}
  \Cov(z_s,z)\,J_{O_s,:}^{\top} &= \Cov(z_s,z_{U_s})\,J_{O_s,U_s}^{\top},\\
  J_{O_s,:}\,B\,J_{O_s,:}^{\top} &= J_{O_s,U_s}\,B_{U_sU_s}\,J_{O_s,U_s}^{\top},
\end{aligned}
\]
exactly. The analysis at $s$ therefore needs only the ensemble at
$U_s\cup\{s\}$ and the compact Jacobian from (a); no row of $J$ is
truncated. The ensemble-transform update uses the same rows through the
observation-space perturbations $Y = X\,J_{O_s,:}^{\top} = X_{U_s}\,J_{O_s,U_s}^{\top}$
and the innovation of the ensemble mean on $O_s$, so the same conclusion
holds for it. With $\Rdec>0$ the rows $J_{O_s,:}$ are non-zero on every site
within $\Rdec$ of $U_s$, and the same identities hold only after replacing
$U_s$ by that larger set. \qed
